%% ------------------------------------------------------------------------------------
\section{Additional analyses performed}
\label{sec:analyses}

Everything in this section was run on the laptop during the audit (18--19 September). Each analysis is a registered experiment (E15--E21, \texttt{experiments/registry/}) whose script reproduces it under provenance; numbers are from those runs where complete and from the identical-seed scratch computation otherwise (marked $\dagger$).

\subsection{The Student-$t$ envelope cell is a handful of exceedances (E11 re-read)}
\label{sec:exceed}
E11's chains were regenerated from their tags. Sites with $|a|>8$ among $8192$ per seed: $1, 1, 1, 2$ (seeds 0--3); the analytic expectation for the variance-normalised $t_5$ at $\rho=0$ is $8192\times P(|t_5|>8/\sqrt{0.6})=1.2$. Seed 2's exceedance is at $|a|=15.4$, beyond the $A=12$ reference; seed 3's at $10.0$. Every ``boundary belief mass $=1.00$'' in the frozen file is one of these sites; the drift is flat in $N_g$ because a pinned site is pinned on every grid; and the pooled relative drift at the lowest level ($0.046$ for seed 2) is one site's score error $\approx(\alpha/\Delta)(15.4-8)$ against the energy of $8192$ ordinary sites. At $\rho=0.5$ and $0.85$ no site exceeds $8$ in any seed, yet the ``boundary belief mass'' reaches $0.09$ and $0.17$: sites at $|a|\approx6$--$7$ whose posterior at the lowest level has a tail beyond $8$. Same mechanism, no pinning, drift $\le10^{-2}$.

The correct diagnostic for polynomial tails is the expected exceedance count $ML\,P(|a|>A)$, Poisson-distributed, and the correct statement is that a finite domain cannot make it zero: $A=12$ gives $0.17$ per $8192$ sites, $A=16$ gives $0.04$. A second mechanism in the same table, mentioned nowhere: the uniform family's worst-level drift of $0.115$ at $\rho=0.5$ on $(401,10)$ is the trapezoidal rule on a discontinuous kernel ($O(h)$, not $O(h^2)$), largest at the coarsest spacing.

\subsection{The closure gap on a wide domain (E18, part ``chain'')}
\label{sec:e18}
$\rho=0.85$, 256 chains $\times$ 2, levels $1$--$5$, three grids. Slopes are paired across grids (same chains) and across levels (common noise) and are precise; ratios to the constant from 512 chains are not (next subsection).

\begin{center}\small
\begin{tabular}{@{}llrrrrrr@{}}
\toprule
family & grid & slope $[3,3.5]$ & $[3.5,4]$ & $[4,4.5]$ & $[4.5,5]$ & ratio $t=4$ & ratio $t=5$ \\
\midrule
Student-$t_5$ & $(801,8)$ & $-4.97$ & $-4.21$ & $-4.04$ & $-4.01$ & 27 & 1440 \\
 & $(1601,16)$ & $-7.69$ & $-7.20$ & $-5.73$ & $-4.48$ & 0.99 & 18 \\
 & $(3201,32)$ & $-7.79$ & $-7.90$ & $-7.86$ & $-7.43$ & 0.74 & 1.06 \\
Student-$t_7$ & $(801,8)$ & $-7.71$ & $-6.86$ & $-5.13$ & $-4.21$ & 1.94 & 54 \\
 & $(3201,32)$ & $-7.90$ & $-7.96$ & $-7.98$ & $-7.99$ & 1.08 & 1.09 \\
Student-$t_9$ & $(1601,16)$ & $-7.98$ & $-8.01$ & $-8.02$ & $-8.02$ & 1.04 & 1.02 \\
Laplace & $(801,8)$ & $-7.95$ & $-8.00$ & $-8.02$ & $-8.02$ & 1.59 & 1.56 \\
bimodal & $(3201,32)$ & $-7.87$ & $-7.95$ & $-7.98$ & $-7.99$ & 1.13 & 1.15 \\
uniform & any & $-7.52$ & $-6.25$ & $-4.62$ & $-4.08$ & 3.8 & 148 \\
\bottomrule
\end{tabular}
\end{center}

Three readings. (i) Every departure from $-8$ on the $(801,8)$ grid for the Student-$t$ families is removed by widening the domain: it was the truncated prior in the reference, not the moment condition. (ii) The uniform family's departure is identical on every grid: it is the double-precision floor of the BP posterior means (a squared gap of $\sim10^{-13}$ per chain), which every family reaches near $t\approx5$; the figure shades it. (iii) On a domain wide enough for the sample, Student-$t_5$ follows $e^{-8t}$ through $t=5$ and reaches the constant. The untruncated \emph{scalar} quadrature (part ``scalar'', adaptive domain $3X/\alpha$) tells the other half: for the unit-variance $t_5$ law the local slope of the expected gap is $-5.9, -6.1, -5.9, -5.6, -5.5, -5.4, -5.3, -5.3$ on consecutive half-unit intervals from $t=1$, converging to $-\nu=-5$, and its ratio to $\tfrac16\kappa_4^2e^{-8t}$ grows as $e^{3t}$ ($0.7, 2.0, 5.0, 14, 48, 170, 630, 2400, 9400$ at $t=1,\dots,5$); for $t_7$ the ratio is $1.0$--$1.4$ through $t=5$ (its crossover is later), for $t_9$ it is $1.00$. The expected tail term is carried by sites with $|a|\gtrsim e^{2t}$: at $t=3$ a site contributing to it has probability $\sim10^{-9}$, and $8192\times32$ sites contain none. Proposition~\ref{prop:heavy} is therefore confirmed in expectation and irrelevant in any sample of practical size; the paper says both.

\subsection{The closure constant, and how precisely it is verified (E18, part ``constant'')}
\label{sec:const}
The per-chain squared gap is a degree-six polynomial in nearly Gaussian variables and is heavy-tailed: over $8192$ chains its maximum is $\approx300$ times its median. A 512-chain mean therefore carries a relative standard error of $\approx15\%$: two 512-chain replicates of E01's Laplace closure at $t=3$ give ratios $0.83$ and $1.14$, two 256-chain replicates $1.77$ and $1.33$. E01's ``$0.99$ at $t=3$'' is the mean of four such replicates and its precision is $\approx\pm10\%$, not $1\%$. With $8192$ chains at $(1601,16)$ and a chain bootstrap:

\begin{center}\small
\begin{tabular}{@{}lcc@{}}
\toprule
family & ratio at $t=3$ & ratio at $t=4$ \\
\midrule
Laplace & 0.92 [0.84, 1.01] & 0.99 [0.91, 1.08] \\
bimodal & 0.93 [0.83, 1.05] & 0.98 [0.88, 1.11] \\
uniform & 0.96 [0.89, 1.04] & 3.6 (floor) \\
Student-$t_5$ & 0.70 [0.62, 0.79] & 1.03 [0.99, 1.09] \\
Student-$t_7$ & 0.90 [0.82, 0.98] & 0.96 [0.88, 1.04] \\
Student-$t_9$ & 0.88 [0.81, 0.96] & 0.93 [0.86, 1.01] \\
\bottomrule
\end{tabular}
\end{center}
Every family with an exponential moment is within its interval of $1$ at $t=4$ (the uniform family has hit the double-precision floor there, ratio $3.6$ with a tiny interval---a floor, not a sample); the heavy-tailed members approach from below and $t_5$ is at $1.03$ at $t=4$. The paper quotes these intervals; this does not touch the rates, which are paired and measured to $\pm0.05$.

\subsection{The locality floor (E15)}
\label{sec:e15}
Pooled over the schedule, the error of the exact best radius-$r$ estimator (sum-product on the window with the exact marginal of its first site):

\begin{center}\small
\begin{tabular}{@{}lrrrrrrrrrr@{}}
\toprule
$\rho$ & $r=0$ & $1$ & $2$ & $3$ & $4$ & $6$ & $8$ & $12$ & $16$ & LMMSE (full) \\
\midrule
0.5 & 0.419 & 0.068 & 0.016 & 0.005 & 0.002 & 0.0003 & 0.0000 & 0.0000 & 0.0000 & 0.264 \\
0.85 & 0.666 & 0.216 & 0.094 & 0.050 & 0.030 & 0.013 & 0.007 & 0.002 & 0.0006 & 0.170 \\
0.95 & 0.755 & 0.317 & 0.174 & 0.107 & 0.072 & 0.039 & 0.025 & 0.011 & 0.006 & 0.081 \\
\bottomrule
\end{tabular}
\end{center}
The floor rises with $\rho$ ($0.002$, $0.030$, $0.072$ at radius $4$), and at $\rho=0.95$ the radius-4 floor is close to the LMMSE floor itself: there, giving up the far sites costs about as much as giving up everything beyond second order. The Gaussian twin's windowed LMMSE reaches its full value by $r=3$ at $\rho=0.5$, so everything a wider window buys is non-Gaussian information. At $\rho=0.85$ the radius-4 floor peaks at $0.065$ around $t\approx0.7$--$1.0$ and is $0.030$ pooled---\emph{below} the window head's $0.118$ (all sites) and $0.075$ (interior) by a factor of two and a half to four, and close to EM--BP's $0.023$. Verdict: locality is \emph{not} the ceiling the manuscript and the report describe. A radius-4 estimator can in principle come within $0.007$ pooled of the exact-inference estimator; the window head's plateau is a property of that one-layer architecture and its training under the recipe, and the manuscript's sentence ``a radius-4 window cannot represent a message that has travelled the length of the chain'' is withdrawn (the message's information beyond radius 4 is worth $0.030$). What survives, sharpened: locality is a good prior at small $M$, and the plateau of the local architecture at large $M$ is an approximation limit that the exact floor now quantifies. The TCN (receptive radius 15) at $0.068$ and the global architectures at $0.060$ are likewise far above their floors ($0.0006$ at $r=16$; $0$ for the full chain), which says the same thing about all of them: at $M=2048$ under this recipe, every network is limited by approximation and optimisation, not by what its receptive field carries.

\subsection{The shape-information constant (E21, part ``constant'')}
\label{sec:scalar-const}
Scalar generalised Gaussian, untruncated quadrature: $I_t(\beta)e^{8t}/[\tfrac1{24}(\partial_\beta\kappa_4)^2]=1.0025, 1.0003, 1.0000$ at $t=1, 1.5, 2$ for $\beta=2$; $1.042, 1.006, 1.001$ for $\beta=3$; $0.960, 0.995, 0.999$ for $\beta=1.5$; the variance coordinate's $\tfrac12$ to five digits. The constant of eq.~(5) is correct. On the chain (E01) the outer-product estimate plateaus at $0.81$ ($\beta=2$) and $0.92$ ($\beta=3$) of the prediction from $t=3$ on, and diverges upward for $\beta\le1.5$ beyond $t\approx2.5$: the estimator's floor (a cancellation of $O(1)$ terms to $\alpha^4\approx10^{-5}$ against $10^{-5}$ grid accuracy) is reached before the asymptote for the kinked members and biases the smooth ones. The manuscript may say ``rate verified on the chain; constant verified in the scalar case and to $20\%$ on the chain''.

\subsection{The fourth-order correction (E16)}
\label{sec:e16}
Laplace chain, fraction of the squared closure gap closed by $\Sg+\nabla\epsilon/(1+\epsilon)$: $-11$ at $t=0.05$ (the expansion diverges; $\epsilon$ reaches $43$), $-1.4$ at $0.15$, $+0.30$ at $0.22$, $0.83$ at $0.32$, $0.92$ at $0.47$, $0.96$ at $0.68$, $0.985$ at $0.98$, $0.996$ at $1.4$, $1.000$ from $2.1$. Bimodal and uniform: the same shape, converging from $t\approx0.15$. Student-$t_5$: $0.13$ at $0.68$, $0.59$ at $1.4$, $0.82$ at $2.1$, then $0.39, 0.08, 0.01$ at $3, 3.5, 4$---on the $(801,8)$ grid, where the reference is truncation-dominated at those levels (Section~\ref{sec:e18}); a rerun at the wide domain is a one-line change and must be made before the number is used. Pooled over the twelve levels the corrected score is \emph{worse} than the twin ($0.47$ vs $0.155$ on Laplace) because the lowest levels dominate the pooled energy and the expansion diverges there: the fourth-order oracle is a per-level statement about what the gap is made of, not a competitor on the schedule. That is the right way to present it.

\subsection{Closed-form misspecification bias (E17)}
\label{sec:e17}
\begin{center}\small
{\small\setlength{\tabcolsep}{4pt}\begin{tabular}{@{}lrrrrrrr@{}}
\toprule
mechanism & $\varepsilon$ & $\rho^\ast$ & EM $\hat\rho$ & $b(\varepsilon)$ & lag-1 plug-in & EM--BP $512/2048$ & pred.\ $512/2048$ \\
\midrule
AR(2) & 0.1 & 0.867 & 0.855 & 0.060 & 0.052 & 0.056 / 0.055 & 0.063 / 0.061 \\
 & 0.4 & 0.916 & 0.888 & 0.248 & 0.216 & 0.214 / 0.214 & 0.248 / 0.248 \\
site-varying & 0.05 & 0.848 & 0.836 & 0.088 & 0.088 & 0.094 / 0.094 & 0.090 / 0.088 \\
 & 0.14 & 0.850 & 0.804 & 0.268 & 0.268 & 0.282 / 0.287 & 0.269 / 0.269 \\
long-range & 0.1 & 0.792 & 0.746 & 0.139 & 0.111 & 0.120 / 0.120 & 0.140 / 0.139 \\
 & 0.4 & 0.713 & 0.659 & 0.235 & 0.202 & 0.206 / 0.204 & 0.235 / 0.235 \\
rank-one & 0.5 & 0.886 & 0.884 & 0.028 & 0.022 & 0.030 / 0.027 & 0.034 / 0.030 \\
 & 1.0 & 0.937 & 0.932 & 0.065 & 0.048 & 0.055 / 0.052 & 0.068 / 0.066 \\
\midrule
\multicolumn{8}{@{}l}{small-$\varepsilon$ laws: AR(2) $b\approx0.60\,\varepsilon$; site-varying $b\approx1.75\,\varepsilon$;}\\
\multicolumn{8}{@{}l}{long-range $b\approx1.07\,\varepsilon^{0.8\to1.0}$; rank-one $b\approx0.15\,\beta^{2.00}$}\\
\bottomrule
\end{tabular}}
\end{center}
The KL projection under the one-view protocol reproduces EM--BP's $M=2048$ error within $0.003$ for site-varying $\rho$ and within $10$--$15\%$ for AR(2) and long-range, where the lag-one plug-in AR(1) (E08's \texttt{lmmse\_ar1} arm) is the closer match: the fitted $\hat\rho$ sits between the two projections because the $C=8$ kernel is not exactly Gaussian and the finite-sample fit under-shoots the KL optimum. Either way the asymptotic bias is a closed-form function of $\varepsilon$ and $L$; its order is $\varepsilon^1$ for the three mechanisms that perturb $O(L)$ directions of the covariance and $\varepsilon^2$ (in $\beta$) for the rank-one latent, which perturbs one. With $v=0.2$ from the well-specified cells, $\sqrt{b^2+v/M}$ predicts the two-size pattern of E08 at every strength.

\subsection{Oracle-free checkpoint selection (E20)}
Eight seeds at $M=32$, the size at which the choice of rung matters most (at $M=2048$ the frozen validation curves are monotone through the ladder, $0.0147\to0.0145$ over the last $400$ iterations, so every rule picks the top). Three seeds in: the oracle rule (validation risk against the true kernel) and the held-out one-view log-evidence rule select the \emph{same} rung in every case (rungs $20$, $40$, $40$), so the test errors are identical ($0.130$, $0.050$, $0.121$), against $0.232$, $0.187$, $0.158$ at the last rung of the ladder. The oracle is not what makes the selection work; the held-out evidence sees the same optimum. \textsc{[five seeds pending; each cell costs about forty minutes because the mixture M-step at $N_g=401$ is five seconds per iteration whatever $M$ is.]}

\subsection{Moment estimators (E19)}
\label{sec:e19}
Same chains and noise draws as E06, $M=512$, sixteen seeds.

\begin{center}\small
\begin{tabular}{@{}lrrrrrrrr@{}}
\toprule
$t$ & 0.05 & 0.1 & 0.2 & 0.4 & 0.7 & 1.0 & 1.4 & 2.0 \\
\midrule
RMSE $\hat\rho$, moments & 0.031 & 0.023 & 0.030 & 0.039 & 0.049 & 0.060 & 0.16 & 0.39 \\
RMSE $\hat\rho$, EM--BP & 0.005 & 0.005 & 0.006 & 0.008 & 0.011 & 0.009 & 0.03 & 0.34 \\
RMSE $\hat\kappa_4$, moments & 1.42 & 1.21 & 1.82 & 1.84 & 3.78 & 15 & 12700 & 1220 \\
RMSE $\hat\kappa_4$, EM--BP & 0.28 & 0.48 & 0.33 & 0.73 & 1.55 & 2.79 & 2.92 & 3.06 \\
\bottomrule
\end{tabular}
\end{center}
 Unbiased (bias $\le0.4$ up to $t=0.4$), never collapsing, $5$--$30\times$ the reference in dispersion at the informative levels, and with the asymptotic growth visible where the estimate is already useless: OLS log-slopes on $t\ge1$ of $1.83$ ($\hat\rho$, predicted $2$) and $3.73$ ($\hat\kappa_4$, predicted $4$). This is the high-noise competitor of the brief's Part~VII.6, and it answers cleanly: rate-correct, inefficient, and the inefficiency is the value of exact inference below the crossover (at high $t$ the MLE's score is itself the degree-4 Hermite statistic, so the two estimators become asymptotically equivalent---one sentence in the paper).

\subsection{The information curve at $\rho=0.5$ (E21, part ``curve'')}
\label{sec:e21}
Table~\ref{tab:crb}. Local slopes of $\log(1/I^{\mathrm{eff}}_\beta)$: $5.5, 5.7, 6.1, 6.5, 6.8, 6.9, 6.6$ on $[0.4,1.2]$; of $\log(1/I^{\mathrm{eff}}_\rho)$: $3.8$--$3.9$ throughout. E12b's levels $[0.6,1.0]$ sit where the shape slope is $6$--$7$. The corollary's exponents are the $t\to\infty$ limits; the experiment measures the curve. Amendment recorded in \texttt{experiments/registry/E12b.json} before any cell of that job ran.

\subsection{Interfaces and paired factors (E03 re-read)}
Table~\ref{tab:matched}. The ISM window head is the matched arm at every size; its paired factor at $M=2048$ is $7.9$ $[6.6,9.3]$; the fixed-corruption DSM arm, which sees the clean chain, is at $7.3$ $[6.2,8.7]$. Ratio-of-means versions: $7.2$ and $6.7$. Both should be given, with the estimand named as in Appendix~C.

\subsection{Compute accounting}
Single-thread timings on the laptop (under load; relative values are what matter), $L=32$:

\begin{center}\small
\begin{tabular}{@{}lrrr@{}}
\toprule
 & $N_g=201$ & $N_g=401$ & $N_g=801$ \\
\midrule
BP inference, one sequence, one level & 0.8 ms & 1.8 ms & 6.2 ms \\
E-step, one sequence & 1.0 ms & 2.5 ms & 9.1 ms \\
mixture M-step ($C=8$), per EM iteration & 1.5 s & 5.2 s & 22 s \\
\midrule
\bottomrule
\end{tabular}\\[6pt]
\begin{tabular}{@{}lrrrrr@{}}
\toprule
 & window & bi-GRU & U-Net & TCN & transformer \\
\midrule
forward pass, one sequence, one level & 0.11 ms & 0.20 ms & 0.22 ms & 0.33 ms & 0.41 ms \\
training, 32k steps (frozen \texttt{wall}) & 113 s & 525 s & 630 s & 641 s & 481 s \\
ISM training, 32k steps & 3{,}400 s$^\dagger$ & --- & --- & --- & 1{,}700 s$^\ddagger$ \\
\bottomrule
\end{tabular}\\[2pt]
{\footnotesize $^\dagger$ exact divergence; $^\ddagger$ Hutchinson estimator.}
\end{center}
So at the operating grid: inference over the twelve-level schedule costs EM--BP $\approx21$ ms per sequence against $1$--$5$ ms for a network (a factor $5$--$20$); an EM--BP fit at $M=2048$ with $\approx1000$ iterations costs $\approx(5+5)\,\mathrm{s}\times1000\approx3$ hours single-threaded against $2$--$11$ minutes for a network (a factor $\approx20$--$100$), with the mixture M-step, not the E-step, dominating below $M\approx2000$; memory is $O(N_g^2)$ per kernel and $O(ML N_g)$ per E-step chunk. The Pareto sentence for the paper: a factor $\approx8$ in error at equal data, bought with one structural assumption (priced by E17) and one to two orders of magnitude of compute. Scaling: linear in $L$ and $M$, quadratic in $N_g$, linear in $C$ for the M-step.
